\documentclass[../../main.tex]{subfiles}

\graphicspath{{../../image/}{../../figure/}}

\begin{document}

\section{凸熵与经典解的存在性}\label{section:ch05-凸熵与经典解的存在性}

本节旨在确立具凸熵 $\eta$ 的齐次守恒律方程组 Cauchy 问题局部经典解的存在性：
\begin{align}
\partial_{t}U(x,t)+\sum_{\alpha=1}^{m}\partial_{\alpha}G_{\alpha}(U(x,t))=0,\qquad x\in\mathbb{R}^{m},\ t>0,\label{align:ch05-cauchy-system}
\end{align}
\begin{align}
U(x,0)=U_{0}(x),\qquad x\in\mathbb{R}^{m},\label{align:ch05-cauchy-data}
\end{align}
其中通量 $G$、熵 $\eta$ 及相应的熵通量 $Q$ 均为定义在 $\mathbb{R}^{n}$ 中以原点为心的闭球 $\overline{\mathscr{B}_{\rho}}$ 上的光滑函数。

贯穿本章，我们采用如下记号。一个\emph{多重指标} $r$ 是指非负整数的 $m$ 元组：$r=(r_{1},\ldots,r_{m})$。记 $|r|=r_{1}+\cdots+r_{m}$ 为 $r$ 的阶数，并记 $\partial^{r}=\partial^{r_{1}}\cdots\partial^{r_{m}}$。于是 $\partial^{r}$ 是一个 $|r|$ 阶微分算子。对 $\ell=-1,0,1,\ldots$，$H_{\ell}$ 与 $H_{\ell}^{m}$ 分别表示 Sobolev 空间 $W^{l,2}(\mathbb{R}^{m};\mathbb{R}^{n})$ 与 $W^{l,2}(\mathbb{R}^{m};\mathbb{M}^{n\times m})$，相应的范数为 $\|\cdot\|_{\ell}$ 与 $\|\|\cdot\|\|_{\ell}$。特别地，$H_{\ell}^{m}$ 等同于 Cartesian 乘积空间 $[H_{\ell}]^{m}$。我们还以符号 $\nabla$ 表示梯度算子 $(\partial_{1},\ldots,\partial_{m})$。于是 $V\in H_{\ell}$ 蕴涵 $\nabla V\in H_{\ell-1}^{m}$，且 $\|\|\nabla V\|\|_{\ell-1}\leqslant\|V\|_{\ell}$。由 Sobolev 嵌入定理，当 $\ell>\frac{m}{2}+1$ 时，$H_{\ell}$ 连续嵌入 $\mathbb{R}^{m}$ 上连续可微 $n$-向量场空间 $C^{1}(\mathbb{R}^{m};\mathbb{R}^{n})$。

对 $U\in\overline{\mathscr{B}_{\rho}}$，我们引入 $n\times n$ 矩阵
\begin{align}
A(U)=\mathrm{D}^{2}\eta(U),\label{align:ch05-A-def}
\end{align}
\begin{align}
J_{\alpha}(U)=A(U)\mathrm{D}G_{\alpha}(U),\qquad\alpha=1,\cdots,m,\label{align:ch05-Jalpha-def}
\end{align}
由\eqref{align:ch03-可积性}，它们是对称的。

本节的主要结果如下。

\begin{theorem}\label{theorem:ch05-凸熵局部存在}
假设守恒律方程组\eqref{align:ch05-cauchy-system}装备有凸熵 $\eta$，从而对任意 $U\in\overline{\mathscr{B}_{\rho}}$，$A(U)$ 正定。设初值 $U_{0}\in H_{\ell}$，其中某 $\ell>\frac{m}{2}+1$，取值于半径 $\rho_{0}<\rho$ 的球 $\overline{\mathscr{B}_{\rho_{0}}}$。则存在 $T_{\infty}\leqslant\infty$ 及 $\mathbb{R}^{m}\times[0,T_{\infty})$ 上唯一连续可微、取值于 $\overline{\mathscr{B}_{\rho}}$ 的函数 $U$，它是 Cauchy 问题\eqref{align:ch05-cauchy-system}、\eqref{align:ch05-cauchy-data} 在时间区间 $[0,T_{\infty})$ 上的经典解。此外，
\begin{align}
U(\cdot,t)\in\bigcap_{k=0}^{\ell}C^{k}([0,T_{\infty});H_{\ell-k}).\label{align:ch05-regularity}
\end{align}
区间 $[0,T_{\infty})$ 是极大的，即若 $T_{\infty}<\infty$，则
\begin{align}
\int_{0}^{T_{\infty}}\|\|\nabla U(\cdot,t)\|\|_{L^{\infty}}\,\mathrm{d}t=\infty\label{align:ch05-blowup1}
\end{align}
且/或
\begin{align}
\limsup_{t\to T_{\infty}}\|U(\cdot,t)\|_{L^{\infty}}=\rho.\label{align:ch05-blowup2}
\end{align}
\end{theorem}

上述定理的传统证明——见\cref{section:ch05-注释}所引文献，甚至在本书第二版中亦然——是在某函数空间 $\mathscr{F}$ 中，将 $V\in\mathscr{F}$ 映为线性化方程组
\begin{align}
\partial_{t}U(x,t)+\sum_{\alpha=1}^{m}\mathrm{D}G_{\alpha}(V(x,t))\partial_{\alpha}U(x,t)=0,\label{align:ch05-linearized}
\end{align}
在初始条件\eqref{align:ch05-cauchy-data}下的解 $U\in\mathscr{F}$ 的映射的不动点，以此确定\eqref{align:ch05-cauchy-system}、\eqref{align:ch05-cauchy-data}的解。
该方法在熵 $\eta$ 凸时有效，因为此时以 $A(V)$ 左乘可使方程组\eqref{align:ch05-linearized}对称化；然而，它在\cref{section:ch05-对合与偶然熵}将遇到的条件下不适用，那里熵不凸，补偿性估计不可避免地与\eqref{align:ch05-cauchy-system}的几何结构相联系，不能照搬到线性化形式\eqref{align:ch05-linearized}。

因此，我们在此采用粘性消去法，即将\eqref{align:ch05-cauchy-system}的解确定为抛物方程组
\begin{align}
\partial_{t}U(x,t)+\sum_{\alpha=1}^{m}\partial_{\alpha}G_{\alpha}(U(x,t))=\varepsilon\Delta U(x,t)\label{align:ch05-parabolic}
\end{align}
的解当 $\varepsilon\to0$ 时的极限。

这一途径未必经由最直接的路线导出\cref{theorem:ch05-凸熵局部存在}的证明，但它具有使本章后续各节的过渡尽可能平滑的优点。粘性消去法的另一好处是它从初等水平出发：唯一的预备知识是会解经典热方程的 Cauchy 问题。

第一步是对固定 $\varepsilon>0$ 建立\eqref{align:ch05-parabolic}、\eqref{align:ch05-cauchy-data} Cauchy 问题的局部存在性。\eqref{align:ch05-parabolic}中的主项是 Laplace 算子，故在此阶段熵不起任何作用。

贯穿本章，我们用 $c$ 表示某一般正常数，它至多可依赖于 $\rho$ 以及 $G$、$\eta$ 及其在 $\overline{\mathscr{B}_{\rho}}$ 上导数的界。

\begin{lemma}\label{lemma:ch05-parabolic-local}
如\cref{theorem:ch05-凸熵局部存在}的陈述所设，假设 $U_{0}$ 取值于 $\overline{\mathscr{B}_{\rho_{0}}}$，$\rho_{0}<\rho$，且属于 $H_{\ell}$，$\ell>\frac{m}{2}+1$。令 $\omega_{0}=\|U_{0}\|_{\ell}$。则对任意固定的 $\omega>\omega_{0}$ 与 $\varepsilon>0$，存在 $T_{\omega,\varepsilon}$，$0<T_{\omega,\varepsilon}\leqslant\infty$，以及\eqref{align:ch05-parabolic}、\eqref{align:ch05-cauchy-data}在时间区间 $[0,T_{\omega,\varepsilon})$ 上的一个解 $U$，取值于 $\overline{\mathscr{B}_{\rho}}$，满足
\begin{align}
U(\cdot,t)\in C^{0}([0,T_{\omega,\varepsilon});H_{\ell})\bigcap L^{2}([0,T_{\omega,\varepsilon});H_{\ell+1}),\label{align:ch05-lemma-spaces}
\end{align}
且
\begin{align}
\|U(\cdot,t)\|_{\ell}<\omega,\qquad 0\leqslant t<T_{\omega,\varepsilon}.\label{align:ch05-lemma-bound}
\end{align}
此外，若 $T_{\omega,\varepsilon}<\infty$，则 $\limsup_{t\to T_{\omega,\varepsilon}}\|U(\cdot,t)\|_{\ell}=\omega$ 且/或 $\limsup_{t\to T_{\omega,\varepsilon}}\|U(\cdot,t)\|_{L^{\infty}}=\rho$。
\end{lemma}

\begin{proof}
固定 $\omega_{1}$，使 $\omega_{0}<\omega_{1}<\omega$。待定的 $T>0$，我们考虑 Lipschitz 函数类 $\mathscr{V}$：$V$ 定义在 $\mathbb{R}^{m}\times[0,T]$ 上，取值于 $\overline{\mathscr{B}_{\rho}}$，满足
\begin{align}
V(\cdot,t)\in L^{\infty}([0,T];H_{\ell}),\qquad\sup_{[0,T]}\|V(\cdot,t)\|_{\ell}\leqslant\omega_{1}.\label{align:ch05-V-class}
\end{align}
由 $L^{p}$ 范数的标准弱下半连续性，$\mathscr{V}$ 在度量
\begin{align}
d(V,\bar{V})=\sup_{[0,T]}\|V(\cdot,t)-\bar{V}(\cdot,t)\|_{0}\label{align:ch05-V-metric}
\end{align}
下是完备度量空间。

对任一给定 $V\in\mathscr{V}$，我们在 $\mathbb{R}^{m}\times[0,T]$ 上构造线性抛物方程组（耦合热方程）
\begin{align}
\partial_{t}U(x,t)-\varepsilon\Delta U(x,t)=-\sum_{\alpha=1}^{m}\partial_{\alpha}G_{\alpha}(V(x,t)),\label{align:ch05-linear-parabolic}
\end{align}
在初始条件\eqref{align:ch05-cauchy-data}下的解 $U$。于是
\begin{align}
(4\pi\varepsilon)^{\frac{m}{2}}U(x,t)&=\int_{\mathbb{R}^{m}}t^{-\frac{m}{2}}\exp\left[-\frac{|x-y|^{2}}{4\varepsilon t}\right]U_{0}(y)\,\mathrm{d}y\nonumber\\
&-\int_{0}^{t}\int_{\mathbb{R}^{m}}(t-\tau)^{-\frac{m}{2}}\exp\left[-\frac{|x-y|^{2}}{4\varepsilon(t-\tau)}\right]\sum_{\alpha=1}^{m}\partial_{\alpha}G_{\alpha}(V(y,\tau))\,\mathrm{d}y\,\mathrm{d}\tau.\label{align:ch05-Duhamel}
\end{align}

我们来建立 $U$ 的先验界，特别地证明当 $T$ 适当小时 $U\in\mathscr{V}$，且映 $V\mapsto U$ 的映射是压缩的。该映射的唯一不动点即为\eqref{align:ch05-parabolic}、\eqref{align:ch05-cauchy-data}在时间区间 $[0,T]$ 上的解。

首先，由\eqref{align:ch05-Duhamel}与\eqref{align:ch05-V-class}，
\begin{align}
\|U(\cdot,t)-U_{0}(\cdot)\|_{L^{\infty}}\leqslant c\omega(\sqrt{\varepsilon t}+t),\label{align:ch05-Linf-bound}
\end{align}
对 $0\leqslant t\leqslant T$ 成立，特别地表明当 $T$ 充分小时，对任意 $t\in[0,T]$，$U(\cdot,t)$ 取值于 $\mathscr{B}_{\rho}$。

固定任意阶 $|r|\leqslant\ell$ 的多重指标 $r$，令 $U_{r}=\partial^{r}U$、$U_{0r}=\partial^{r}U_{0}$，对\eqref{align:ch05-linear-parabolic}作用 $\partial^{r}$，得
\begin{align}
\partial_{t}U_{r}(x,t)-\varepsilon\Delta U_{r}(x,t)=-\sum_{\alpha=1}^{m}\partial_{\alpha}\partial^{r}G_{\alpha}(V(x,t)).\label{align:ch05-Ur-eq}
\end{align}
由于 $\|V(\cdot,t)\|_{L^{\infty}}<\rho$ 且 $\|V(\cdot,t)\|_{\ell}<\omega$，Sobolev 空间理论中熟知的插值估计给出
\begin{align}
\|\partial^{r}G_{\alpha}(V(\cdot,t))\|_{0}\leqslant c\|V(\cdot,t)\|_{\ell}\leqslant c\omega,\label{align:ch05-Gr-interp}
\end{align}
对任意 $|r|\leqslant\ell$ 与任意 $t\in[0,T]$ 成立。于是由热方程的标准理论，$U_{r}$ 作为以 $U_{0r}$ 为初值的\eqref{align:ch05-Ur-eq}的解，属于 $C^{0}([0,T];H_{0})\bigcap L^{2}([0,T];H_{1})$。这给出两个"能量"积分，即
\begin{align}
\int_{\mathbb{R}^{m}}|U_{r}(x,t)|^{2}\,\mathrm{d}x+2\varepsilon\int_{0}^{t}\int_{\mathbb{R}^{m}}|\nabla U_{r}|^{2}\,\mathrm{d}x\,\mathrm{d}\tau
&=\int_{\mathbb{R}^{m}}|U_{0r}(x)|^{2}\,\mathrm{d}x+2\int_{0}^{t}\int_{\mathbb{R}^{m}}\sum_{\alpha=1}^{m}\partial_{\alpha}U_{r}^{\top}\partial^{r}G_{\alpha}(V)\,\mathrm{d}x\,\mathrm{d}\tau\nonumber\\
&\leqslant\int_{\mathbb{R}^{m}}|U_{0r}(x)|^{2}\,\mathrm{d}x+\varepsilon\int_{0}^{t}\int_{\mathbb{R}^{m}}|\nabla U_{r}|^{2}\,\mathrm{d}x\,\mathrm{d}\tau+\frac{ct\omega^{2}}{\varepsilon},\label{align:ch05-energy1}
\end{align}
它由以 $2U_{r}^{\top}$ 乘\eqref{align:ch05-Ur-eq}、在 $\mathbb{R}^{m}\times[0,t]$ 上积分并分部积分形式导出，对任意 $|r|\leqslant\ell$ 成立；以及
\begin{align}
2\int_{0}^{t}\int_{\mathbb{R}^{m}}|\partial_{\tau}U_{r}|^{2}\,\mathrm{d}x\,\mathrm{d}\tau+\varepsilon\int_{\mathbb{R}^{m}}|\nabla U_{r}(x,t)|^{2}\,\mathrm{d}x
&=\varepsilon\int_{\mathbb{R}^{m}}|\nabla U_{0r}(x)|^{2}\,\mathrm{d}x-2\int_{0}^{t}\int_{\mathbb{R}^{m}}\partial_{\tau}U_{r}^{\top}\sum_{\alpha=1}^{m}\partial_{\alpha}\partial^{r}G_{\alpha}(V)\,\mathrm{d}x\,\mathrm{d}\tau\nonumber\\
&\leqslant\int_{0}^{t}\int_{\mathbb{R}^{m}}|\partial_{\tau}U_{r}|^{2}\,\mathrm{d}x\,\mathrm{d}\tau+c(t+\varepsilon)\omega^{2},\label{align:ch05-energy2}
\end{align}
它由以 $2\partial_{t}U_{r}^{\top}$ 乘\eqref{align:ch05-Ur-eq}、在 $\mathbb{R}^{m}\times[0,t]$ 上积分并分部积分形式导出，对任意 $|r|\leqslant\ell-1$ 成立。

对所有 $|r|\leqslant\ell$ 的 $r$ 求和\eqref{align:ch05-energy1}，对所有 $|r|\leqslant\ell-1$ 的 $r$ 求和\eqref{align:ch05-energy2}，我们导出估计
\begin{align}
\|U(\cdot,t)\|_{\ell}^{2}\leqslant\omega_{0}^{2}+\frac{ct\omega^{2}}{\varepsilon},\label{align:ch05-Uell-bound}
\end{align}
\begin{align}
\int_{0}^{t}\|\partial_{\tau}U(\cdot,\tau)\|_{\ell-1}^{2}\,\mathrm{d}\tau\leqslant(\varepsilon+ct)\omega^{2}.\label{align:ch05-dtU-bound}
\end{align}
特别地，\eqref{align:ch05-Uell-bound}蕴涵当 $T$ 充分小时 $\sup_{[0,T]}\|U(\cdot,t)\|_{\ell}<\omega_{1}$，从而 $U\in\mathscr{V}$。

现固定 $\mathscr{V}$ 中的 $V$ 与 $\bar{V}$，考虑由它们诱导的\eqref{align:ch05-linear-parabolic}、\eqref{align:ch05-cauchy-data}的解 $U$ 与 $\bar{U}$。则
\begin{align}
\partial_{t}(U-\bar{U})-\varepsilon\Delta(U-\bar{U})=-\sum_{\alpha=1}^{m}\partial_{\alpha}[G_{\alpha}(V)-G_{\alpha}(\bar{V})].\label{align:ch05-difference-eq}
\end{align}
以 $2(U-\bar{U})^{\top}$ 乘\eqref{align:ch05-difference-eq}，在 $\mathbb{R}^{m}\times[0,t]$、$t\in(0,T)$ 上积分并分部积分，我们推得
\begin{align}
\int_{\mathbb{R}^{m}}|U(x,t)-\bar{U}(x,t)|^{2}\,\mathrm{d}x+2\varepsilon\int_{0}^{t}\int_{\mathbb{R}^{m}}|\nabla(U-\bar{U})|^{2}\,\mathrm{d}x\,\mathrm{d}\tau
&=2\int_{0}^{t}\int_{\mathbb{R}^{m}}\sum_{\alpha=1}^{m}\partial_{\alpha}(U-\bar{U})^{\top}[G_{\alpha}(V)-G_{\alpha}(\bar{V})]\,\mathrm{d}x\,\mathrm{d}\tau\nonumber\\
&\leqslant\varepsilon\int_{0}^{t}\int_{\mathbb{R}^{m}}|\nabla(U-\bar{U})|^{2}\,\mathrm{d}x\,\mathrm{d}\tau+\frac{cT}{\varepsilon}\sup_{[0,T]}\int_{\mathbb{R}^{m}}|V(x,\tau)-\bar{V}(x,\tau)|^{2}\,\mathrm{d}x.\label{align:ch05-contraction}
\end{align}
回忆\eqref{align:ch05-V-metric}，我们结论：若 $\frac{cT}{\varepsilon}=\mu^{2}<1$，则 $d(U,\bar{U})\leqslant\mu d(V,\bar{V})$，这就证明了映射 $V\mapsto U$ 有唯一不动点，即\eqref{align:ch05-parabolic}、\eqref{align:ch05-cauchy-data}在时间区间 $[0,T]$ 上的唯一解 $U$。

令 $\rho_{1}=\|U(\cdot,T)\|_{0}\leqslant\rho$。若 $\rho_{1}=\rho$，则引理已证毕，取 $T_{\omega,\varepsilon}=T$。另一方面，若 $\rho_{1}<\rho$ 且由于 $\|U(\cdot,T)\|_{\ell}\leqslant\omega_{1}<\omega$，我们可按上述方式以新的初值 $U(\cdot,T)$ 及 $(\rho_{1},\omega_{1})$ 取代 $(\rho_{0},\omega_{0})$ 来解一个新的 Cauchy 问题，将\eqref{align:ch05-parabolic}、\eqref{align:ch05-cauchy-data}的解 $U$ 延拓到 $T$ 之后。迭代此过程，将 $U$ 延拓到极大时间区间 $[0,T_{\omega,\varepsilon})$，要么 $T_{\omega,\varepsilon}=\infty$，要么 $T_{\omega,\varepsilon}<\infty$，此时 $\limsup_{t\to T_{\omega,\varepsilon}}\|U(\cdot,t)\|_{L^{\infty}}=\rho$ 且/或 $\limsup_{t\to T_{\omega,\varepsilon}}\|U(\cdot,t)\|_{\ell}=\omega$。\cref{lemma:ch05-parabolic-local}的证明完成。
\end{proof}

下一步是导出\eqref{align:ch05-parabolic}、\eqref{align:ch05-cauchy-data}的解在 $\varepsilon\to0$ 时仍一致成立的界。注意\eqref{align:ch05-Linf-bound}与\eqref{align:ch05-dtU-bound}对 $t\in[0,T_{\omega,\varepsilon})$ 成立并满足此要求。下述命题确立了该方向上的主要估计，此时凸熵走到前台，粘性项退居辅助地位。

\begin{lemma}\label{lemma:ch05-uniform-bound}
假设由\eqref{align:ch05-A-def}定义的 $A(U)$ 对 $U\in\overline{\mathscr{B}_{\rho}}$ 正定。则存在 $c_{0}>1$，仅依赖于 $\rho$、$\ell$ 以及 $G$、$\eta$ 及其在 $\overline{\mathscr{B}_{\rho}}$ 上导数的界，具有如下性质。设 $U_{0}\in H_{\ell}$，$\ell>\frac{m}{2}+1$，取值于 $\overline{\mathscr{B}_{\rho_{0}}}$，$\rho_{0}<\rho$。固定任意 $\omega\geqslant c_{0}\|U_{0}\|_{\ell}$，并参照\cref{lemma:ch05-parabolic-local}，考虑满足\eqref{align:ch05-lemma-bound}的\eqref{align:ch05-parabolic}、\eqref{align:ch05-cauchy-data}解 $U$ 的寿命 $T_{\omega,\varepsilon}$。则 $T_{\omega,\varepsilon}>T_{\omega}$，其中 $T_{\omega}$ 是不依赖于 $\varepsilon\in(0,1)$ 的正常数，且对任意 $t\in[0,T_{\omega}]$，
\begin{align}
\|U(\cdot,t)\|_{\ell}\leqslant c_{0}\|U_{0}(\cdot)\|_{\ell}\exp\int_{0}^{t}g(\tau)\,\mathrm{d}\tau,\label{align:ch05-uniform-est}
\end{align}
其中
\begin{align}
g(t)\leqslant c\left[\|\|\nabla U(\cdot,t)\|\|_{L^{\infty}}+\varepsilon\|\|\nabla U(\cdot,t)\|\|_{L^{\infty}}^{2}\right].\label{align:ch05-g-def}
\end{align}
\end{lemma}

\begin{proof}
对任意阶 $|r|\leqslant\ell$ 的多重指标 $r$，我们在\eqref{align:ch05-Ur-eq}中取 $V=U$，写成
\begin{align}
\partial_{t}U_{r}+\sum_{\alpha=1}^{m}\mathrm{D}G_{\alpha}(U)\partial_{\alpha}U_{r}
=\sum_{\alpha=1}^{m}\left\{\mathrm{D}G_{\alpha}(U)\partial_{\alpha}U_{r}-\partial^{r}[\mathrm{D}G_{\alpha}(U)\partial_{\alpha}U]\right\}+\varepsilon\Delta U_{r}.\label{align:ch05-Ur-eq2}
\end{align}
在上式右端的求和中，最高阶 $l+1$ 阶导数项相消。因此，由 Sobolev 空间中熟知的插值估计（通常称为 Moser 估计）——见\cref{section:ch05-注释}所引文献——可得
\begin{align}
\|\mathrm{D}G_{\alpha}(U)\partial_{\alpha}U_{r}-\partial^{r}[\mathrm{D}G_{\alpha}(U)\partial_{\alpha}U]\|_{0}\leqslant c\|\|\nabla U\|\|_{L^{\infty}}\|\|\nabla U\|\|_{\ell-1}.\label{align:ch05-Moser}
\end{align}
于是\eqref{align:ch05-Ur-eq2}诱导出如下"能量"积分：
\begin{align}
\int_{\mathbb{R}^{m}}U_{r}^{\top}(x,t)A(U(x,t))U_{r}(x,t)\,\mathrm{d}x-\int_{\mathbb{R}^{m}}U_{0r}^{\top}(x)A(U_{0}(x))U_{0r}(x)\,\mathrm{d}x
&=-2\varepsilon\int_{0}^{t}\int_{\mathbb{R}^{m}}\partial_{\alpha}U_{r}^{\top}A(U)\partial_{\alpha}U_{r}\,\mathrm{d}x\,\mathrm{d}\tau-4\varepsilon\int_{0}^{t}\int_{\mathbb{R}^{m}}\sum_{\alpha=1}^{m}U_{r}^{\top}\partial_{\alpha}A(U)\partial_{\alpha}U_{r}\,\mathrm{d}x\,\mathrm{d}\tau\nonumber\\
&\quad-\varepsilon\int_{0}^{t}\int_{\mathbb{R}^{m}}\sum_{\alpha=1}^{m}U_{r}^{\top}[\partial_{\alpha}\mathrm{D}A(U)\partial_{\alpha}U]U_{r}\,\mathrm{d}x\,\mathrm{d}\tau\nonumber\\
&\quad+\int_{0}^{t}\int_{\mathbb{R}^{m}}\sum_{\alpha=1}^{m}U_{r}^{\top}[\partial_{\alpha}J_{\alpha}(U)-\mathrm{D}A(U)\partial_{\alpha}G_{\alpha}(U)]U_{r}\,\mathrm{d}x\,\mathrm{d}\tau\nonumber\\
&\quad+2\int_{0}^{t}\int_{\mathbb{R}^{m}}\sum_{\alpha=1}^{m}U_{r}^{\top}A(U)[\mathrm{D}G_{\alpha}(U)\partial_{\alpha}U_{r}-\partial^{r}[\mathrm{D}G_{\alpha}(U)\partial_{\alpha}U]]\,\mathrm{d}x\,\mathrm{d}\tau,\label{align:ch05-energy-A}
\end{align}
它由以 $2U_{r}^{\top}A(U)$ 乘\eqref{align:ch05-Ur-eq2}、在 $\mathbb{R}^{m}\times[0,t]$ 上积分并分部积分形式导出。在此过程中使用了如下恒等式，它们体现了矩阵 $A(U)$ 与 $J_{\alpha}(U)$ 的对称性：
\begin{align}
2U_{r}^{\top}A(U)\partial_{t}U_{r}=\partial_{t}[U_{r}^{\top}A(U)U_{r}]-U_{r}^{\top}\partial_{t}A(U)U_{r},\label{align:ch05-id1}
\end{align}
\begin{align}
U_{r}^{\top}\partial_{t}A(U)U_{r}=-U_{r}^{\top}\left[\sum_{\alpha=1}^{m}\mathrm{D}A(U)\partial_{\alpha}G_{\alpha}(U)\right]U_{r}+\varepsilon U_{r}^{\top}[\mathrm{D}A(U)\Delta U]U_{r},\label{align:ch05-id2}
\end{align}
\begin{align}
U_{r}^{\top}[\mathrm{D}A(U)\Delta U]U_{r}&=\sum_{\alpha=1}^{m}\partial_{\alpha}[U_{r}^{\top}[\partial_{\alpha}A(U)]U_{r}]\nonumber\\
&\quad-2\sum_{\alpha=1}^{m}U_{r}^{\top}[\partial_{\alpha}A(U)]\partial_{\alpha}U_{r}-\sum_{\alpha=1}^{m}U_{r}^{\top}[\partial_{\alpha}\mathrm{D}A(U)\partial_{\alpha}U]U_{r},\label{align:ch05-id3}
\end{align}
\begin{align}
U_{r}^{\top}A(U)\Delta U_{r}&=\sum_{\alpha=1}^{m}\partial_{\alpha}[U_{r}^{\top}A(U)\partial_{\alpha}U_{r}]\nonumber\\
&\quad-\sum_{\alpha=1}^{m}\partial_{\alpha}U_{r}^{\top}A(U)\partial_{\alpha}U_{r}-\sum_{\alpha=1}^{m}U_{r}^{\top}\partial_{\alpha}A(U)\partial_{\alpha}U_{r},\label{align:ch05-id4}
\end{align}
\begin{align}
\sum_{\alpha=1}^{m}2U_{r}^{\top}A(U)\mathrm{D}G_{\alpha}(U)\partial_{\alpha}U_{r}
=\sum_{\alpha=1}^{m}\partial_{\alpha}[U_{r}^{\top}J_{\alpha}(U)U_{r}]-\sum_{\alpha=1}^{m}U_{r}^{\top}\partial_{\alpha}J_{\alpha}(U)U_{r}.\label{align:ch05-id5}
\end{align}
由于 $A(U)$ 正定，
\begin{align}
\int_{\mathbb{R}^{m}}\partial_{\alpha}U_{r}^{\top}(x,\tau)A(U(x,\tau))\partial_{\alpha}U_{r}(x,\tau)\,\mathrm{d}x\geqslant\mu\int_{\mathbb{R}^{m}}|\partial_{\alpha}U_{r}(x,\tau)|^{2}\,\mathrm{d}x.\label{align:ch05-A-positive}
\end{align}
由 Cauchy--Schwarz 不等式，
\begin{align}
2\int_{\mathbb{R}^{m}}U_{r}^{\top}(x,\tau)\partial_{\alpha}A(U(x,\tau))\partial_{\alpha}U_{r}(x,\tau)\,\mathrm{d}x
\leqslant\mu\int_{\mathbb{R}^{m}}|\partial_{\alpha}U_{r}(x,\tau)|^{2}\,\mathrm{d}x+c\int_{\mathbb{R}^{m}}|\partial_{\alpha}U(x,\tau)|^{2}|U_{r}(x,\tau)|^{2}\,\mathrm{d}x.\label{align:ch05-CS}
\end{align}
因此，对所有 $|r|\leqslant\ell$ 求和\eqref{align:ch05-energy-A}，并利用\eqref{align:ch05-A-positive}、\eqref{align:ch05-CS}与\eqref{align:ch05-Moser}，我们得到
\begin{align}
\sum_{|r|\leqslant\ell}\int_{\mathbb{R}^{m}}U_{r}^{\top}(x,t)A(U(x,t))U_{r}(x,t)\,\mathrm{d}x
\leqslant\sum_{|r|\leqslant\ell}\int_{\mathbb{R}^{m}}U_{0r}^{\top}(x)A(U_{0}(x))U_{0r}(x)\,\mathrm{d}x+c\int_{0}^{t}g(\tau)\|U(\cdot,\tau)\|_{\ell}^{2}\,\mathrm{d}\tau,\label{align:ch05-sum-energy}
\end{align}
其中 $g$ 按\eqref{align:ch05-g-def}有界。又因 $A(U)$ 正定，
\begin{align}
\sum_{|r|\leqslant\ell}\int_{\mathbb{R}^{m}}U_{r}^{\top}(x,t)A(U(x,t))U_{r}(x,t)\,\mathrm{d}x\geqslant\mu\|U(\cdot,t)\|_{\ell}^{2}.\label{align:ch05-sum-positive}
\end{align}
因此，\eqref{align:ch05-sum-energy}给出
\begin{align}
\|U(\cdot,t)\|_{\ell}^{2}\leqslant c_{0}^{2}\|U_{0}(\cdot)\|_{\ell}^{2}+2\int_{0}^{t}g(\tau)\|U(\cdot,\tau)\|_{\ell}^{2}\,\mathrm{d}\tau,\label{align:ch05-Gronwall-setup}
\end{align}
由此由 Gronwall 引理得\eqref{align:ch05-uniform-est}。

由于 $\|\|\nabla U(\cdot,\tau)\|\|_{L^{\infty}}\leqslant c\|U(\cdot,\tau)\|_{\ell}\leqslant c\omega$，由\eqref{align:ch05-uniform-est}显然存在 $T_{\omega}>0$，使得 $c_{0}\|U_{0}\|_{\ell}<\omega$ 蕴涵对所有 $t\in[0,T_{\omega}]$ 有 $\|U(\cdot,t)\|_{\ell}<\omega$，即对所有 $\varepsilon\in(0,1)$ 有 $T_{\omega}<T_{\omega,\varepsilon}$。引理证明完成。
\end{proof}

至此，我们已通过在\eqref{align:ch05-parabolic}中令 $\varepsilon\to0$，为处理双曲方程组\eqref{align:ch05-cauchy-system}搭好了舞台。

\begin{lemma}\label{lemma:ch05-limit}
假设 $A(U)$ 对 $U\in\overline{\mathscr{B}_{\rho}}$ 正定，并取 $U_{0}\in H_{\ell}$，$\ell>\frac{m}{2}+1$，取值于 $\overline{\mathscr{B}_{\rho_{0}}}$，$\rho_{0}<\rho$。固定 $\omega\geqslant c_{0}\|U_{0}\|_{\ell}$，其中 $c_{0}$ 是\cref{lemma:ch05-uniform-bound}中引入的常数，并确定相应的 $T_{\omega}$。则存在\eqref{align:ch05-cauchy-system}、\eqref{align:ch05-cauchy-data}的一个经典解 $U$，定义在 $\mathbb{R}^{m}\times[0,T_{\omega}]$ 上，取值于 $\overline{\mathscr{B}_{\rho}}$。对任意 $t\in[0,T_{\omega}]$，$U(\cdot,t)\in H_{\ell}$，$\|U(\cdot,t)\|_{\ell}\leqslant\omega$，且
\begin{align}
\|U(\cdot,t)\|_{\ell}\leqslant c_{0}\|U_{0}(\cdot)\|_{\ell}\exp\int_{0}^{t}c\|\|\nabla U(\cdot,\tau)\|\|_{L^{\infty}}\,\mathrm{d}\tau.\label{align:ch05-limit-est}
\end{align}
此外，函数 $t\mapsto U(\cdot,t)$ 在 $[0,T_{\omega}]$ 上按 $H_{\ell}$ 连续。
\end{lemma}

\begin{proof}
取任意序列 $\{\varepsilon_{k}\}$，$\varepsilon_{k}\to0$（$k\to\infty$），令 $U_{k}$ 为\eqref{align:ch05-parabolic}、\eqref{align:ch05-cauchy-data}在 $\varepsilon=\varepsilon_{k}$、时间区间 $[0,T_{\omega}]$ 上的解。由\cref{lemma:ch05-uniform-bound}，$\{U_{k}\}$ 在 $L^{\infty}([0,T_{\omega}];H_{\ell})$ 中有界，且对 $0\leqslant t\leqslant T_{\omega}$ 有 $\|U_{k}(\cdot,t)\|_{\ell}<\omega$。此外，由\eqref{align:ch05-dtU-bound}知 $\{U_{k}\}$ 在 $W^{1,2}([0,T_{\omega}];H_{\ell-1})$ 中亦有界。因此，由 Sobolev 空间的标准理论，$\{U_{k}\}$ 等度连续，从而包含一子列（仍记为 $\{U_{k}\}$），在 $\mathbb{R}^{m}\times[0,T_{\omega}]$ 的紧集上一致收敛到某连续函数 $U$，取值于 $\overline{\mathscr{B}_{\rho}}$，在分布意义下满足\eqref{align:ch05-cauchy-system}、\eqref{align:ch05-cauchy-data}。对任意 $t\in[0,T_{\omega}]$，界 $\|U_{k}(\cdot,t)\|_{\ell}<\omega$ 蕴涵 $U_{k}(\cdot,t)\rightharpoonup U(\cdot,t)$ 弱收敛于 $H_{\ell}$，且 $\|U(\cdot,t)\|_{\ell}\leqslant\omega$。这又给出 $U(\cdot,t)\in C^{1}(\mathbb{R}^{m})$，$\|\|\nabla U(\cdot,t)\|\|_{L^{\infty}}\leqslant c\omega$；且因 $U$ 是\eqref{align:ch05-cauchy-system}的解，$\partial_{t}U(\cdot,t)\in C^{0}(\mathbb{R}^{m})$，$\|\partial_{t}U(\cdot,t)\|_{L^{\infty}}\leqslant c\omega$。故 $U$ 是 Lipschitz 的。

接下来我们证明对任意 $t\in[0,T_{\omega}]$，当 $k\to\infty$ 时 $U_{k}(\cdot,t)\to U(\cdot,t)$ 在 $H_{0}$ 中强收敛。为此注意到 $V_{k}=U_{k}-U$ 满足方程
\begin{align}
\partial_{t}V_{k}+\sum_{\alpha=1}^{m}\mathrm{D}G_{\alpha}(U)\partial_{\alpha}V_{k}=\varepsilon\Delta U_{k}-\sum_{\alpha=1}^{m}[\mathrm{D}G_{\alpha}(U_{k})-\mathrm{D}G_{\alpha}(U)]\partial_{\alpha}U_{k},\label{align:ch05-Vk-eq}
\end{align}
在 $\mathbb{R}^{m}\times[0,T_{\omega}]$ 上，初始条件 $V_{k}(\cdot,0)=0$。于是以 $2V_{k}^{\top}A(U)$ 乘\eqref{align:ch05-Vk-eq}，在 $\mathbb{R}^{m}\times[0,t]$、$t\in[0,T_{\omega}]$ 上积分并分部积分，得
\begin{align}
\int_{\mathbb{R}^{m}}V_{k}^{\top}(x,t)A(U(x,t))V_{k}(x,t)\,\mathrm{d}x
&=\int_{0}^{t}\int_{\mathbb{R}^{m}}V_{k}^{\top}\{\partial_{t}A(U)+\sum_{\alpha=1}^{m}\partial_{\alpha}J_{\alpha}(U)\}V_{k}\,\mathrm{d}x\,\mathrm{d}\tau\nonumber\\
&\quad-2\int_{0}^{t}\int_{\mathbb{R}^{m}}V_{k}^{\top}A(U)[\mathrm{D}G_{\alpha}(U_{k})-\mathrm{D}G_{\alpha}(U)]\partial_{\alpha}U_{k}\,\mathrm{d}x\,\mathrm{d}\tau\nonumber\\
&\quad-2\varepsilon_{k}\int_{0}^{t}\int_{\mathbb{R}^{m}}\sum_{\alpha=1}^{m}\partial_{\alpha}[V_{k}^{\top}A(U)]\partial_{\alpha}U_{k}\,\mathrm{d}x\,\mathrm{d}\tau.\label{align:ch05-Vk-energy}
\end{align}
由于 $A(U)$ 正定，
\begin{align}
\int_{\mathbb{R}^{m}}V_{k}^{\top}(x,t)A(U(x,t))V_{k}(x,t)\,\mathrm{d}x\geqslant\mu\|V_{k}(\cdot,t)\|_{0}^{2}.\label{align:ch05-Vk-positive}
\end{align}
因此，\eqref{align:ch05-Vk-energy}诱导出形如
\begin{align}
\|V_{k}(\cdot,t)\|_{0}^{2}\leqslant c\omega\int_{0}^{t}\|V_{k}(\cdot,\tau)\|_{0}^{2}\,\mathrm{d}\tau+c\omega^{2}\varepsilon_{k}^{2}t\label{align:ch05-Vk-Gronwall}
\end{align}
的不等式，由此我们结论：对任意 $t\in[0,T_{\omega}]$，当 $k\to\infty$ 时 $V_{k}(\cdot,t)\to0$ 在 $H_{0}$ 中强收敛。

由于 $\{U_{k}(\cdot,t)\}$ 在 $H_{0}$ 中强收敛到 $U(\cdot,t)$、在 $H_{\ell}$ 中弱收敛，由插值知收敛在 $H_{\ell-1}$ 中是强的，且在 $\mathbb{R}^{m}$ 上一致。此外，对任意 $t\in[0,T_{\omega}]$，$\nabla U_{k}(\cdot,t)\to\nabla U(\cdot,t)$ 在 $\mathbb{R}^{m}$ 上一致。特别地，回忆\eqref{align:ch05-uniform-est}对 $U_{k}$ 成立、$g$ 按\eqref{align:ch05-g-def}有界，令 $k\to\infty$，即验证 $U$ 满足\eqref{align:ch05-limit-est}。

还需证明 $t\mapsto U(\cdot,t)$ 在 $[0,T_{\omega}]$ 上按 $H_{\ell}$ 连续。考虑到\eqref{align:ch05-cauchy-system}在时间平移与反射下不变，只需证明 $t\mapsto U(\cdot,t)$ 在 $t=0$ 处右连续，即 $t\to0$ 时 $U(\cdot,t)\to U_{0}(\cdot)$ 在 $H_{\ell}$ 中。

从恒等式
\begin{align}
\sum_{|r|\leqslant\ell}\int_{\mathbb{R}^{m}}U_{r}^{\top}A(U)U_{r}\,\mathrm{d}x-\sum_{|r|\leqslant\ell}\int_{\mathbb{R}^{m}}U_{kr}^{\top}A(U_{k})U_{kr}\,\mathrm{d}x
&=-\sum_{|r|\leqslant\ell}\int_{\mathbb{R}^{m}}U_{kr}^{\top}[A(U_{k})-A(U)]U_{kr}\,\mathrm{d}x\nonumber\\
&\quad-\sum_{|r|\leqslant\ell}\int_{\mathbb{R}^{m}}2(U_{kr}-U_{r})^{\top}A(U)U_{r}\,\mathrm{d}x\nonumber\\
&\quad-\sum_{|r|\leqslant\ell}\int_{\mathbb{R}^{m}}(U_{kr}-U_{r})^{\top}A(U)(U_{kr}-U_{r})\,\mathrm{d}x,\label{align:ch05-continuity-id}
\end{align}
出发，它对任意固定的 $t\in[0,T_{\omega}]$ 成立，令 $k\to\infty$。右端前两项趋于零，而最后一项有确定符号：
\begin{align}
\sum_{|r|\leqslant\ell}\int_{\mathbb{R}^{m}}(U_{kr}-U_{r})^{\top}A(U)(U_{kr}-U_{r})\,\mathrm{d}x\geqslant\mu\|U_{k}-U\|_{\ell}^{2}.\label{align:ch05-continuity-sign}
\end{align}
由此，当 $k\to\infty$ 时\eqref{align:ch05-continuity-id}左端的上极限非正。回忆 $U_{k}$ 满足\eqref{align:ch05-sum-energy}，我们结论
\begin{align}
\sum_{|r|\leqslant\ell}\int_{\mathbb{R}^{m}}U_{r}^{\top}(x,t)A(U(x,t))U_{r}(x,t)\,\mathrm{d}x-\sum_{|r|\leqslant\ell}\int_{\mathbb{R}^{m}}U_{0r}^{\top}(x)A(U_{0}(x))U_{0r}(x)\,\mathrm{d}x\leqslant c\omega^{3}t.\label{align:ch05-continuity-est}
\end{align}

现写出恒等式
\begin{align}
\sum_{|r|\leqslant\ell}\int_{\mathbb{R}^{m}}[U_{r}(x,t)-U_{0r}(x)]^{\top}A(U_{0}(x))[U_{r}(x,t)-U_{0r}(x)]\,\mathrm{d}x
&=\sum_{|r|\leqslant\ell}\int_{\mathbb{R}^{m}}U_{r}^{\top}(x,t)A(U(x,t))U_{r}(x,t)\,\mathrm{d}x\nonumber\\
&\quad-\sum_{|r|\leqslant\ell}\int_{\mathbb{R}^{m}}U_{0r}^{\top}(x)A(U_{0}(x))U_{0r}(x)\,\mathrm{d}x\nonumber\\
&\quad-\sum_{|r|\leqslant\ell}\int_{\mathbb{R}^{m}}U_{r}^{\top}(x,t)[A(U(x,t))-A(U_{0}(x))]U_{r}(x,t)\,\mathrm{d}x\nonumber\\
&\quad-\sum_{|r|\leqslant\ell}\int_{\mathbb{R}^{m}}2[U_{r}(x,t)-U_{0r}(x)]^{\top}A(U_{0}(x))U_{0r}(x)\,\mathrm{d}x,\label{align:ch05-continuity-id2}
\end{align}
令 $t\to0$。右端最后一项趋于零，因为 $U$ 是连续函数且 $\|U(\cdot,t)\|_{\ell}\leqslant\omega$，由此知 $t\mapsto U(\cdot,t)$ 至少在 $H_{\ell}$ 中弱连续。类似地，倒数第二项趋于零，因为 $U_{k}$ 满足\eqref{align:ch05-Linf-bound}，且当 $k\to\infty$ 时该估计传递到 $U$。最后，由\eqref{align:ch05-continuity-est}，剩余两项的贡献非正。于是我们结论：当 $t\to0$ 时\eqref{align:ch05-continuity-id2}左端的上极限非正。另一方面，由于 $A(U)$ 正定，
\begin{align}
\sum_{|r|\leqslant\ell}\int_{\mathbb{R}^{m}}[U_{r}(x,t)-U_{0r}(x)]^{\top}A(U_{0}(x))[U_{r}(x,t)-U_{0r}(x)]\,\mathrm{d}x\geqslant\mu\|U(\cdot,t)-U_{0}(\cdot)\|_{\ell}^{2}.\label{align:ch05-continuity-positive}
\end{align}
因此，$t\mapsto U(\cdot,t)$ 在 $H_{\ell}$ 中连续。特别地，$U$ 是\eqref{align:ch05-cauchy-system}、\eqref{align:ch05-cauchy-data}的连续可微经典解。由\eqref{align:ch05-Linf-bound}，必要时可缩小 $T_{\omega}$ 以保证对 $t\in[0,T_{\omega}]$ 有 $\|U(\cdot,t)\|_{L^{\infty}}<\rho$。引理证明完成。
\end{proof}

\emph{\cref{theorem:ch05-凸熵局部存在}的证明。}我们从\cref{lemma:ch05-limit}在 $[0,T_{\omega}]$ 上构造的\eqref{align:ch05-cauchy-system}、\eqref{align:ch05-cauchy-data}的解 $U$ 出发，置 $\rho_{1}=\|U(\cdot,T_{\omega})\|_{\infty}<\rho$、$\omega_{1}=\|U(\cdot,T_{\omega})\|_{\ell}$。于是我们可以初值 $U(\cdot,T_{\omega})$ 及 $(\rho_{1},c_{0}\omega_{1})$ 取代 $(\rho_{0},\omega)$ 来解一个新的\eqref{align:ch05-cauchy-system} Cauchy 问题，将 $U$ 延拓到 $T_{\omega}$ 之后。迭代此过程，将 $U$ 延拓到极大时间区间 $[0,T_{\infty})$，要么 $T_{\infty}=\infty$，要么 $T_{\infty}<\infty$，此时 $\limsup_{t\to T_{\infty}}\|U(\cdot,t)\|_{L^{\infty}}=\rho$ 或 $\limsup_{t\to T_{\infty}}\|U(\cdot,t)\|_{\ell}=\infty$。由\eqref{align:ch05-limit-est}，$\|U(\cdot,t)\|_{\ell}$ 仅当\eqref{align:ch05-blowup1}成立时才可能无界。

$U$ 的时间导数与混合时空导数可通过空间导数借助方程组\eqref{align:ch05-cauchy-system}确定。于是由已建立的 $U(\cdot,t)\in C^{0}([0,T_{\infty});H_{\ell})$ 即得\eqref{align:ch05-regularity}。

解的唯一性将在\cref{section:ch05-相对熵与经典解的稳定性}相当弱的正则性假设下建立。

如第 IV 章所指出，$T_{\infty}=\infty$ 是罕见情形：一般地，光滑解在有限时间内激波发展时破裂。

\cref{theorem:ch05-凸熵局部存在}证明的关键在于对称正定矩阵值函数 $A(U)$ 的存在，它通过使\eqref{align:ch05-Jalpha-def}定义的矩阵值函数 $J_{\alpha}(U)$ 对称而充当\emph{对称化子}（symmetrizer）。此处 $A(U)$ 恰为 $\eta(U)$ 的 Hessian 矩阵纯属巧合，因为存在装备有并非来自熵的对称化子的方程组。事实上，即使对装备有所谓\emph{symbolic symmetrizers}的方程组，Cauchy 问题解的存在性也已建立。特别地，这包括所有特征速度 $\lambda_{i}(v;U)$ 的重数不随 $v$ 或 $U$ 变化的双曲方程组。

\end{document}
